% Part: first-order-logic
% Chapter: natural-deduction
% Section: derivations

\documentclass[../../../include/open-logic-section]{subfiles}

\begin{document}

\iftag{FOL}
      {\olfileid{fol}{ntd}{der}}
      {\olfileid{pl}{ntd}{der}}

\olsection{\usetoken{P}{derivation}}

\begin{explain}
Kita wis nerangake apa kang diarani asumsi, lan wis menehi aturan
inferensi. !!^{derivation}s ing deduksi natural dibangkitake kanthi
induktif saka perkara-perkara kasebut: saben !!{derivation} dadi
asumsi dhewe, utawa dumadi saka siji, loro, utawa telung
!!{derivation}s kang diterusake dening inferensi kang bener.
\end{explain}

\begin{defn}[!!^{derivation}]
\Article{derivation} \emph{!!{derivation}} saka !!a{sentence}~$!A$
adhedhasar asumsi~$\Gamma$ yaiku wit winates saka !!{sentence}s kang
nyukupi syarat-syarat iki:
\begin{enumerate}
\item !!{sentence}s paling dhuwur ing wit kalebu ing $\Gamma$ utawa
  !!{discharged} dening sawijining inferensi ing wit.
\item !!{sentence} paling ngisor ing wit yaiku~$!A$.
\item Saben !!{sentence} ing wit kajaba ukara~$!A$ ing sisih ngisor
  iku premis saka panrapan aturan inferensi kang bener, lan
  kesimpulan aturan iku mapan langsung ing sangisore !!{sentence}
  kasebut ing wit.
\end{enumerate}
Banjur kita ngarani $!A$ \emph{kesimpulan} saka !!{derivation}
kasebut lan $\Gamma$ minangka asumsi kang !!{undischarged}.

Yen ana !!a{derivation} saka $!A$ adhedhasar~$\Gamma$, kita kandha yen
$!A$ \emph{!!{derivable}} saka~$\Gamma$, utawa kanthi simbol:
$\Gamma \Proves !A$. Yen ana !!a{derivation} saka~$!A$ kang saben
asumsine wis !!{discharged}, kita nulis~$\Proves !A$.
\end{defn}

\begin{ex}
Saben asumsi dhewe iku !!a{derivation}. Dadi, tuladhane, $!A$ dhewe
iku !!a{derivation}, lan $!B$ dhewe uga. Kita bisa ngasilake
!!{derivation} anyar saka loro iku kanthi ngetrapake, umpamane, aturan
$\Intro{\land}$,
\begin{prooftree}
\AxiomC{$!A$}
\AxiomC{$!B$}
\RightLabel{\Intro{\land}}
\BinaryInfC{$!A \land !B$}
\end{prooftree}
Aturan kasebut dimaksudake umum: $!A$ lan~$!B$ ing aturan iku bisa
diganti nganggo sembarang !!{sentence}s, tuladhane $!C$ lan~$!D$.
Banjur kesimpulane dadi $!C \land !D$, mula
\begin{prooftree}
\AxiomC{$!C$}
\AxiomC{$!D$}
\RightLabel{\Intro{\land}}
\BinaryInfC{$!C \land !D$}
\end{prooftree}
iku !!{derivation} kang bener. Mesthi wae, kita uga bisa ngijolake
asumsine, saengga $!D$ dadi~$!A$ lan $!C$ dadi~$!B$. Dadi,
\begin{prooftree}
\AxiomC{$!D$}
\AxiomC{$!C$}
\RightLabel{\Intro{\land}}
\BinaryInfC{$!D \land !C$}
\end{prooftree}
uga !!{derivation} kang bener.

Saiki kita bisa ngetrapake aturan liyane, umpamane $\Intro{\lif}$,
kang ngidini kita nyimpulake kondisional lan ngidini kita
!!{discharge} sembarang asumsi kang padha karo anteseden kondisional
kasebut. Mula, loro
!!{derivation}s iki padha-padha bener:
\begin{prooftree}
\AxiomC{$\Discharge{!C}{1}$}
\AxiomC{$!D$}
\RightLabel{\Intro{\land}}
\BinaryInfC{$!C \land !D$}
\DischargeRule{\Intro{\lif}}{1}
\UnaryInfC{$!C \lif (!C \land !D)$}
\DisplayProof\bottomAlignProof
\AxiomC{$!C$}
\AxiomC{$\Discharge{!D}{1}$}
\RightLabel{\Intro{\land}}
\BinaryInfC{$!C \land !D$}
\DischargeRule{\Intro{\lif}}{1}
\UnaryInfC{$!D \lif (!C \land !D)$}
\end{prooftree}
Loro iku, kanthi urut, nuduhake yen $!D \Proves !C \lif (!C \land !D)$
lan $!C \Proves !D \lif (!C \land !D)$.

Elinga yen mbubarake asumsi iku idin, dudu kuwajiban: kita ora kudu
mbubarake asumsine. Mligine, kita bisa ngetrapake aturan sanajan
asumsi kasebut ora ana ing !!{derivation}. Tuladha iki sah, sanajan
ora ana asumsi~$!A$ kang arep !!{discharged}:
\begin{prooftree}
  \AxiomC{$!B$}
  \DischargeRule{\Intro{\lif}}{1}
  \UnaryInfC{$!A \lif !B$}
\end{prooftree}
\end{ex}

\end{document}
